\section{Q-Chain finality: the 23-axis binding}
\label{sec:qchain-finality}

\subsection{What Q-Chain binds}

Q-Chain is the chain that produces finality certificates over every other
chain in the Lux ecosystem. HIP-0079 defines Q-Chain as the
``post-quantum finality plane'': it consumes EpochCommitments from every
execution chain (C-Chain, P-Chain, X-Chain, Z-Chain), folds them into a
canonical QBlock, and produces a Pulsar finality certificate over the
QBlock digest.

The QBlock digest is the single 48-byte commitment that says ``this is the
state of the Lux Network at epoch $N$ under LUX\_STRICT\_E2E\_PQ.''
Everything else --- the bridge admission, the cross-chain Warp envelope, the
Z-Chain commitment binding into a \texttt{TxAuthEnvelope} verification ---
ultimately points back to this digest.

\subsection{The 23 axes}

\texttt{ComputeRoundDigest} (\texttt{luxfi/consensus@12097034},
\texttt{protocol/qchain/round\_digest.go}) binds 23 axes into the QBlock
hash. The axes are the security-relevant inputs to consensus at this round:

\begin{table}[H]
\centering
\caption{ComputeRoundDigest axes. The 23 axes are the inputs to consensus
at a single round. Adding an axis bumps the customization
\texttt{LUX\_QBLOCK\_V1 $\to$ V2} and is a hard fork.}
\label{tab:qblock-axes}
\begin{tabular}{rlp{8cm}}
\toprule
\# & Axis & Role \\
\midrule
1  & \texttt{ProfileHash}             & ChainSecurityProfile digest. \\
2  & \texttt{Version}                 & QBlock version. \\
3  & \texttt{NetworkID}               & destination network. \\
4  & \texttt{Epoch}                   & monotonic epoch counter. \\
5  & \texttt{Height}                  & Q-Chain height. \\
6  & \texttt{PrevBlockHash}           & parent QBlock digest. \\
7  & \texttt{CChainCommitment}        & C-Chain EpochCommitment digest. \\
8  & \texttt{PChainCommitment}        & P-Chain EpochCommitment digest. \\
9  & \texttt{XChainCommitment}        & X-Chain EpochCommitment digest. \\
10 & \texttt{ZChainCommitment}        & Z-Chain EpochCommitment digest (\texttt{ZRegistryRoots.Hash()}). \\
11 & \texttt{ValidatorSetRoot}        & active validator set Merkle root. \\
12 & \texttt{ValidatorStakeRoot}      & per-validator stake commitment. \\
13 & \texttt{BridgeAdmissionRoot}     & bridge admission state root. \\
14 & \texttt{TxAuthBatchRoot}         & accepted \texttt{TxAuthEnvelope} root for epoch. \\
15 & \texttt{PermitBatchRoot}         & accepted \texttt{PQPermit} root for epoch. \\
16 & \texttt{ZIdentityRoot}           & shortcut into \texttt{ZRegistryRoots.IdentityRoot}. \\
17 & \texttt{P3qProofRoot}            & p3q STARK proof aggregator root for this epoch. \\
18 & \texttt{RandomnessSeed}          & DRBG seed for this epoch. \\
19 & \texttt{PrevFinalityCert}        & previous epoch's Pulsar cert digest. \\
20 & \texttt{Timestamp}               & uint64 epoch timestamp (UTC seconds). \\
21 & \texttt{LineageRoot}             & lineage chain digest (HIP-0079 §5). \\
22 & \texttt{GovernanceRoot}          & active governance state root. \\
23 & \texttt{ExtraHash}               & 48-byte HIP-future-reserved digest. \\
\bottomrule
\end{tabular}
\end{table}

Every axis except \texttt{ExtraHash} is consumed by some downstream
verifier. \texttt{ExtraHash} is reserved for future HIP additions: future
proposals will bind new axes into the \texttt{ExtraHash} preimage so the
QBlock layout itself stays stable across the migration window.

\subsection{ComputeRoundDigest}

The construction:

\begin{lstlisting}[language=Go,caption={\texttt{ComputeRoundDigest} for
QBlock. Source: \texttt{luxfi/consensus@12097034},
\texttt{protocol/qchain/round\_digest.go}.}]
func (q *QBlock) ComputeRoundDigest() [48]byte {
    parts := [][]byte{
        []byte("Lux/QBlock/v1"),
        q.ProfileHash[:],
        u16BE(q.Version),
        u32BE(q.NetworkID),
        u64BE(q.Epoch),
        u64BE(q.Height),
        q.PrevBlockHash[:],
        q.CChainCommitment[:], q.PChainCommitment[:],
        q.XChainCommitment[:], q.ZChainCommitment[:],
        q.ValidatorSetRoot[:], q.ValidatorStakeRoot[:],
        q.BridgeAdmissionRoot[:],
        q.TxAuthBatchRoot[:], q.PermitBatchRoot[:],
        q.ZIdentityRoot[:], q.P3qProofRoot[:],
        q.RandomnessSeed[:],
        q.PrevFinalityCert[:],
        u64BE(q.Timestamp),
        q.LineageRoot[:],
        q.GovernanceRoot[:],
        q.ExtraHash[:],
    }
    return tupleHash48(parts, "LUX_QBLOCK_V1")
}
\end{lstlisting}

The Pulsar finality lane (Section~\ref{sec:pulsar}) signs over this
digest. The signature, together with the QBlock body, is the finality
certificate.

\subsection{Canonical QBlock layout}

The QBlock body is the wire-serialised form of the struct above plus the
finality certificate. The wire layout is deterministic big-endian, fixed-width
except for the certificate trailer:

\begin{table}[H]
\centering
\caption{QBlock wire layout. The fixed header is 1{,}080 bytes; the
certificate trailer is the Pulsar aggregate plus committee bitmap, sized
per the active committee (typically $\approx 33$ KB).}
\begin{tabular}{lll}
\toprule
Offset & Size & Field \\
\midrule
0       & 48 & \texttt{profile\_hash} \\
48      & 2  & \texttt{version} \\
50      & 4  & \texttt{network\_id} \\
54      & 8  & \texttt{epoch} \\
62      & 8  & \texttt{height} \\
70      & 48 & \texttt{prev\_block\_hash} \\
118     & 48 & \texttt{cchain\_commitment} \\
$\vdots$ & $\vdots$ & $\vdots$ (per Table~\ref{tab:qblock-axes}) \\
1{,}072 & 8  & \texttt{timestamp} \\
1{,}080 & varies & \texttt{finality\_cert} (Pulsar-65 aggregate + bitmap) \\
\bottomrule
\end{tabular}
\end{table}

\subsection{Finality cert binding to Pulsar}

The Pulsar finality lane (Section~\ref{sec:pulsar}) operates as a
two-round threshold signature scheme over the 43-of-64 active committee.
Round 1 commitments hash the QBlock digest as their session id. Round 2
opens partial signatures that aggregate linearly into a single 33-KB
certificate. The verifier:

\begin{enumerate}[leftmargin=2em,nosep]
  \item Recomputes the QBlock digest from the body.
  \item Looks up the active committee against
    \texttt{ValidatorSetRoot}.
  \item Verifies the Pulsar aggregate signature over the digest under the
    committee's group public key (which is itself anchored in
    \texttt{ZRegistryRoots.IdentityRoot} from the Z-Chain).
  \item Confirms the committee bitmap has at least 43 set bits.
\end{enumerate}

Step 2 binds the certificate's committee to the consensus state \emph{at
this epoch}. Step 3 verifies the lattice math. Step 4 enforces the
threshold. The certificate is binding under MLWE.

\subsection{Why 23 axes specifically}

The 23-axis count is the minimum closure of the dependency graph between
the Lux chains. Each axis is the only commitment from \emph{one} subsystem
the verifier needs to confirm that subsystem's state matches what the QBlock
claims. Dropping any axis would force a downstream verifier to consult the
subsystem's own state instead of the QBlock, breaking the ``QBlock is the
single binding'' invariant.

Adding axes is forward-compatible via the \texttt{ExtraHash} reservation;
the next HIP that adds axes binds them into the \texttt{ExtraHash} preimage
and ships an \texttt{ExtraHashV2} struct with the customization
\texttt{LUX\_QBLOCK\_EXTRA\_V2}. The QBlock layout itself does not change
until the migration window closes and a hard-fork to \texttt{LUX\_QBLOCK\_V2}
absorbs the new fields directly.
